% ═══════════════════════════════════════════════════════════════
% AB-02d: Arbeitsblatt Wiederholung Mi familia
% Spanisch Klasse 7 - Sequenz 2
% ═══════════════════════════════════════════════════════════════
% Inhalt: Familienmitglieder, tener, Possessivbegleiter, Text
% Punkte: 20
% ═══════════════════════════════════════════════════════════════

\documentclass[11pt, a4paper]{article}

\newif\ifswmodus
\swmodusfalse

\usepackage[utf8]{inputenc}
\usepackage[T1]{fontenc}
\usepackage[ngerman]{babel}
\usepackage{helvet}
\renewcommand{\familydefault}{\sfdefault}

\usepackage[margin=2cm]{geometry}
\usepackage{enumitem}
\usepackage{tabularx}
\usepackage{booktabs}
\usepackage{multirow}

\usepackage{xcolor}
\usepackage{framed}
\usepackage{tcolorbox}
\tcbuselibrary{skins}

\usepackage{fancyhdr}
\usepackage{lastpage}
\usepackage{needspace}

% FARBEN
\definecolor{spanisch-color}{HTML}{c41e3a}
\definecolor{grau-color}{HTML}{888888}
\definecolor{hellgrau-color}{HTML}{f5f5f5}
\definecolor{orange-color}{HTML}{b35c00}

\definecolor{sw-dark}{HTML}{000000}
\definecolor{sw-gray}{HTML}{666666}
\definecolor{sw-white}{HTML}{ffffff}

\ifswmodus
    \colorlet{spanisch}{sw-dark}
    \colorlet{grau}{sw-gray}
    \colorlet{hellgrau}{sw-white}
    \colorlet{orange}{sw-dark}
\else
    \colorlet{spanisch}{spanisch-color}
    \colorlet{grau}{grau-color}
    \colorlet{hellgrau}{hellgrau-color}
    \colorlet{orange}{orange-color}
\fi

\newcommand{\fachfarbe}{spanisch}

\newcommand{\thema}{Wiederholung: Mi familia}
\newcommand{\sequenz}{02}
\newcommand{\block}{d}
\newcommand{\fach}{Spanisch}
\newcommand{\klasse}{7}
\newcommand{\maxpunkte}{20}
\newcommand{\datum}{\today}

\pagestyle{fancy}
\fancyhf{}
\renewcommand{\headrulewidth}{0.4pt}
\renewcommand{\footrulewidth}{0.4pt}

\lhead{\fach{} -- Klasse \klasse}
\chead{AB-\sequenz\block}
\rhead{\datum}

\lfoot{}
\cfoot{}
\rfoot{Seite \thepage\ von \pageref{LastPage}}

\newcommand{\punktebox}[1]{\hfill\fbox{\small \_/#1}}
\newcommand{\luecke}[1]{\rule{#1}{0.4pt}}

\newtcolorbox{merkkasten}[1][]{
    colback=hellgrau,
    colframe=\fachfarbe,
    colbacktitle=white,
    coltitle=\fachfarbe,
    boxrule=1pt,
    arc=0pt,
    left=8pt, right=8pt, top=8pt, bottom=8pt,
    fonttitle=\bfseries,
    attach boxed title to top left={yshift=-2mm, xshift=5mm},
    boxed title style={boxrule=0pt, colback=white},
    title=#1
}

\newtcolorbox{vokabelkasten}{
    colback=white,
    colframe=\fachfarbe,
    boxrule=0.5pt,
    arc=0pt,
    left=5pt, right=5pt, top=5pt, bottom=5pt
}

\newtcolorbox{hinweis}{
    colback=white,
    colframe=orange,
    colbacktitle=white,
    coltitle=orange,
    boxrule=1pt,
    arc=0pt,
    left=5pt, right=5pt, top=5pt, bottom=5pt,
    fonttitle=\bfseries,
    attach boxed title to top left={yshift=-2mm, xshift=5mm},
    boxed title style={boxrule=0pt, colback=white},
    title=Hinweis
}

\newenvironment{aufgabe}[1]{%
    \needspace{5\baselineskip}%
    \nopagebreak[4]%
    \vspace{10pt}
    \noindent\textbf{\textcolor{\fachfarbe}{Aufgabe #1}}
    \nopagebreak[4]%
    \vspace{5pt}
    \nopagebreak[4]%
    \begin{list}{}{\leftmargin=0pt}
    \item[]
}{%
    \end{list}
}

\newcommand{\es}[1]{\textcolor{\fachfarbe}{\textbf{#1}}}

\begin{document}

% ═══════════════════════════════════════════════════════════════
% KOPFBEREICH
% ═══════════════════════════════════════════════════════════════

\begin{center}
    {\Large\bfseries\textcolor{\fachfarbe}{Arbeitsblatt: \thema}}
\end{center}

\vspace{5pt}

\noindent
\begin{tabularx}{\textwidth}{|X|X|X|}
\hline
\textbf{Name:} \luecke{4cm} &
\textbf{Klasse:} \klasse &
\textbf{Datum:} \luecke{2cm} \\
\hline
\end{tabularx}

\vspace{15pt}

% ═══════════════════════════════════════════════════════════════
% MERKKASTEN: Familienwortschatz
% ═══════════════════════════════════════════════════════════════

\begin{merkkasten}[Merkkasten: La familia]
\textbf{Familienmitglieder:}

\vspace{0.5em}

\begin{tabular}{ll}
\es{la madre} & = \luecke{4cm} \\[0.8em]
\es{el padre} & = \luecke{4cm} \\[0.8em]
\es{la hermana} & = \luecke{4cm} \\[0.8em]
\es{el hermano} & = \luecke{4cm} \\[0.8em]
\es{la abuela} & = \luecke{4cm} \\[0.8em]
\es{el abuelo} & = \luecke{4cm} \\
\end{tabular}

\vspace{0.5em}

\textit{Tipp: Übertrage die Übersetzungen aus dem Lernpfad!}
\end{merkkasten}

\vspace{10pt}

% ═══════════════════════════════════════════════════════════════
% AUFGABE 1: Familienmitglieder zuordnen
% ═══════════════════════════════════════════════════════════════

\begin{aufgabe}{1 -- Familienmitglieder zuordnen}
Verbinde die spanischen Familienbegriffe mit der deutschen Bedeutung. \punktebox{4}
\end{aufgabe}

\begin{center}
\begin{tabular}{rcl}
\es{la madre} & \luecke{3cm} & der Großvater \\[0.8em]
\es{el padre} & \luecke{3cm} & die Mutter \\[0.8em]
\es{la hermana} & \luecke{3cm} & der Vater \\[0.8em]
\es{el abuelo} & \luecke{3cm} & die Schwester \\
\end{tabular}
\end{center}

\vspace{15pt}

% ═══════════════════════════════════════════════════════════════
% AUFGABE 2: tener konjugieren
% ═══════════════════════════════════════════════════════════════

\begin{aufgabe}{2 -- Das Verb \es{tener} (haben)}
Ergänze die richtige Form von \es{tener}. \punktebox{4}
\end{aufgabe}

\begin{vokabelkasten}
\textit{Formen:} tengo | tienes | tiene | tenemos
\end{vokabelkasten}

\vspace{0.5em}

\noindent
a) Yo \luecke{3cm} una hermana. \hfill (ich habe)

\vspace{0.8em}

\noindent
b) Tú \luecke{3cm} dos hermanos. \hfill (du hast)

\vspace{0.8em}

\noindent
c) Ella \luecke{3cm} un perro. \hfill (sie hat)

\vspace{0.8em}

\noindent
d) Nosotros \luecke{3cm} muchos primos. \hfill (wir haben)

\vspace{15pt}

% ═══════════════════════════════════════════════════════════════
% AUFGABE 3: Possessivbegleiter einsetzen
% ═══════════════════════════════════════════════════════════════

\begin{aufgabe}{3 -- Possessivbegleiter einsetzen}
Ergänze den passenden Possessivbegleiter: \es{mi}, \es{tu}, \es{su}, \es{nuestra}, \es{mis}, \es{sus}. \punktebox{6}
\end{aufgabe}

\begin{vokabelkasten}
\textit{Zur Erinnerung:} mi/tu/su = Singular | mis/tus/sus = Plural | nuestro/a = unser/e
\end{vokabelkasten}

\vspace{0.5em}

\noindent
a) \luecke{2cm} madre se llama María. \hfill (meine Mutter)

\vspace{0.8em}

\noindent
b) ¿Cómo se llama \luecke{2cm} hermano? \hfill (dein Bruder)

\vspace{0.8em}

\noindent
c) \luecke{2cm} padre es profesor. \hfill (sein/ihr Vater)

\vspace{0.8em}

\noindent
d) \luecke{2cm} abuela tiene 70 años. \hfill (unsere Großmutter)

\vspace{0.8em}

\noindent
e) \luecke{2cm} hermanos son simpáticos. \hfill (meine Geschwister)

\vspace{0.8em}

\noindent
f) \luecke{2cm} primos viven en Madrid. \hfill (seine/ihre Cousins)

\vspace{15pt}

% ═══════════════════════════════════════════════════════════════
% AUFGABE 4: Text über Familie schreiben
% ═══════════════════════════════════════════════════════════════

\begin{aufgabe}{4 -- Über deine Familie schreiben}
Schreibe einen kurzen Text über deine Familie (mind. 4 Sätze). \punktebox{6}
\end{aufgabe}

\begin{hinweis}
Verwende: Familienmitglieder, \es{tener}, Possessivbegleiter, Alter (Zahlen).

Beispiel: \es{Mi familia es grande. Tengo una madre...}
\end{hinweis}

\vspace{0.5em}

\noindent\rule{\textwidth}{0.4pt}\vspace{8pt}
\noindent\rule{\textwidth}{0.4pt}\vspace{8pt}
\noindent\rule{\textwidth}{0.4pt}\vspace{8pt}
\noindent\rule{\textwidth}{0.4pt}\vspace{8pt}
\noindent\rule{\textwidth}{0.4pt}\vspace{8pt}
\noindent\rule{\textwidth}{0.4pt}

% ═══════════════════════════════════════════════════════════════
% PUNKTEÜBERSICHT
% ═══════════════════════════════════════════════════════════════

\vspace{20pt}
\hrule
\vspace{10pt}

\begin{flushright}
\textbf{Punkte:} \luecke{1cm} / \maxpunkte
\end{flushright}

\end{document}
